\section{Experiments}
\label{sec:exper}


\textbf{Datasets}.
To make state-of-the-art comparison,
We conduct zero-shot evaluation on the popular video benchmarks,
\textit{i.e.},
UCF-101 \cite{Soomro2012UCF101AD}, 
Kinetics-400 \cite{Kay2017TheKH}, 
and 
MSR-VTT \cite{Chen2021TheMT}.
UCF-101 contains videos of 101 action categories.
Following \cite{unterthiner2019accurate},
we use FVD to evaluate the distance between the generated video and the real video.
Kinetics-400 is a dataset comprising videos of 400 action categories. 
Following \cite{li2017video, balaji2019Conditional, wu2022nwa, villegas2023phenaki},
we use FID \cite{heusel2017gans} to assess the performance of video generation.
MSR-VTT is a video dataset with open-vocabulary captions,  
wherein CLIPSIM \cite{Wu2021GODIVAGO} and CLIP-FID \cite{parmar2022aliased} are commonly employed for evaluating the T2V generation.
Moreover,
since these existing benchmarks either have a small number of testing videos or contain only action labels without complex descriptions,
we propose to collect an evaluation benchmark for ablation studies.
It is called as Vimeo11k,
where we collect 11,293 open-world videos along with their captions from 10 mainstream categories in Vimeo.
To our best knowledge,
it is the largest testing benchmark for zero-shot video generation. 
More details and experiments can be found in the supplementary doc.






\noindent\textbf{Implementation Details}. %% 你得有个逻辑，不能想到哪写到哪
\textbf{(1) Director \& Script \& Actor \& Voicer. }
We choose GPT-4 \cite{OpenAI2023GPT4TR} as our LLM director to generate script.
The specific instructions can be found in the supplementary doc.
Then we employ Stable Diffusion XL \cite{Podell2023SDXLIL} and Bark \cite{sunoai2023bark} as our designer and voicer to generate reference actor images and convert scripts into speech.
\textbf{(2) ShowMaker. }
We choose SD-1.4 \cite{rombach2022highResolution} as our base model, 
and follow \cite{singer2023makeavideo} to add temporal self attention.
Then,
we add our temporal text cross attention on top of each temporal self attention.
We expand the input channel of the conv-in layer of U-Net from 4 to 9, 
so that model can take the concatenated feature in Eq. (\ref{eq:ms3}) as input.
We use zero initialization for newly added channels.
We use CLIP ViT-L/14 \cite{radford2021learning} as text encoder $\varepsilon_{T}$,
VQVAE \cite{Esser2020TamingTF} as autoencoder consisting of $\varepsilon$ and $\mathcal{D}$.
Besides,
we use OpenCLIP ViT-H/14 \cite{ilharco2021openclip} as image encoder $\varepsilon_{I}$,
and add spatial image cross attention in the STEB block.
The diffusion step $T$ is set to 1000 as \cite{rombach2022highResolution}.
For training, 
we set the parameter of mode selection as $\alpha=0.6$ and $m=6$ in Eq. (\ref{eq:alpham}),
and $\beta$ is set to 1.
Following \cite{Ho2022ImagenVH, Wang2023LAVIEHV, Ho2022VideoDM, Dandi2019JointlyTI, villegas2023phenaki},
we employ publicly available image dataset (\textit{i.e.}, Laion400M \cite{Schuhmann2021LAION400MOD}) and video dataset (\textit{i.e.}, WebVid10M \cite{bain2021frozen}) for joint training.
More training details can be found in the supplementary doc.



\begin{figure}[t]
    \centering
    \includegraphics[width=0.9\linewidth
    ]{figs/ucf_1024.pdf}
   \vspace{-0.25cm}
    \caption{
    \textbf{Long video generation.}
    Comparison on UCF-101 (1000 frames). 
    The lower FVD, the better generation performance.
    }
    \label{fig:ucf_1024}
    \vspace{-0.3cm}
\end{figure}


\begin{table}[t]
    % \scriptsize
    \centering
    \setlength\tabcolsep{4.5pt}
    \resizebox{1\linewidth}{!}{%
    \begin{tabular}{l c c c}
        \toprule
        \textbf{Method} & \textbf{CLIP-I ($\uparrow$)} & \textbf{CLIP-T ($\uparrow$)} & \textbf{Vlog Duration ($\uparrow$)}  \\
        \midrule
        w/o \small{autoregressive}~\cite{Zhu2023MovieFactoryAM} & 0.5683 & 0.2752 & 1min03s \\
        only \small{autoregressive}~\cite{villegas2023phenaki} & 0.5642 & 0.2535 & 3min58s \\
        \textbf{Ours} & \textbf{0.6294} & \textbf{0.2756} & \textbf{3min58s} \\
        \bottomrule
        \vspace{-1.75em}
    \end{tabular}
    }
    \caption{\textbf{Ablation for generation process}.
    \cite{Zhu2023MovieFactoryAM} cannot generate a long video in a single scene and both \cite{Zhu2023MovieFactoryAM, villegas2023phenaki} can't refer to images.
    }
    \vspace{-1em}
    \label{tab:framework_ablation}
\end{table}



\subsection{Comparison with state-of-the-art}


%\textbf{UCF-101}.
Tab. \ref{tab:ucf} shows that,
no matter whether the input text is the class label or hand-crafted prompt, 
our method achieves the best FVD performance of zero-shot video generation on UCF-101.
%\textbf{Kinetics-400}.
Tab. \ref{tab:k400} shows that,
compared to Phenaki \cite{villegas2023phenaki},
our method achieves a better FID performance of zero-shot setting on Kinetics-400, 
but only using 66.7\% training videos.
Furthermore, our generated videos have a resolution of $320$$\times$$512$, which is higher than Phenaki's $256$$\times$$256$.
%\textbf{MSR-VTT}.
Tab. \ref{tab:msrvtt} shows that,
our method achieves remarkably competitive FID and CLIPSIM performance on MSR-VTT.
Furthermore, 
Fig. \ref{fig:ucf_1024} illustrates that,
we significantly surpass TATS \cite{ge2022long} (\textit{i.e.}, the only open-source long video generation model within our knowledge),
for generating 1000-frame videos on UCF-101.
Moreover,
our method does not encounter the issue of TATS,
where the video quality continuously declines as the number of frames increases.
It is noteworthy that,
our method achieves this performance by zero-shot generation,
without any finetuning on UCF-101.



\subsection{Ablation Study}

\textbf{Vlog Generation Process}.
Tab. \ref{tab:framework_ablation} presents a comparison between different generation processes. 
For a fair comparison, 
we use the same scripts and employ our ShowMaker as videographer. 
We compare three approaches including
no autoregression at all as MovieFactory \cite{Zhu2023MovieFactoryAM}, 
fully autoregressive processes in one go as Phenaki \cite{villegas2023phenaki}, 
and 
the generation process of our Vlogger.
% \textcolor{red}{
Specifically,
We use GPT-4 to generate five stories and go through the Vlogger’s planning process to get five vlog scripts.
We set 20 different random seeds,
so that each of the above generation processes generates 20 different vlogs for each script.
We use CLIP-I, CLIP-T \cite{Ye2023IPAdapterTC} and video duration to evaluate the quality of the generated vlogs. 
The results demonstrate that,
even using the same script and videographer, 
our Vlogger outshines other existing frameworks for preferable vlog generation. 
More details can be found in the supplementary doc.

\begin{figure}[t]
    \centering
    % \vspace{-0.3cm}
    \includegraphics[width=0.9\linewidth
    ]{figs/ip_scale.pdf}
   \vspace{-0.25cm}
    \caption{
    \textbf{Ablation for spatial image cross attention.} 
    The values around a point on the curve correspond to the $\beta$.
    While $\beta$$=$$0$ means the original network without spatial image cross attention.
    }
    \label{fig:ica}
    \vspace{-0.3cm}
\end{figure}






\begin{figure*}[t]
    \centering
    % \vspace{-0.3cm}
    \includegraphics[width=1\textwidth]{figs/movie_case.pdf}
%   \vspace{-0.6cm}
    \caption{
    \textbf{Qualitative comparison with the state-of-the-art methods on long video generation.}
    The story and long video of Phenaki \cite{villegas2023phenaki} are available at
    \href{https://phenaki.github.io/}{phenaki.github.io}. 
    The scene diversity and picture quality generated by Vlogger are significantly better than Phenaki.
    }
    \label{fig:movie_case}
%   \vspace{-0.3cm}
\end{figure*}

\begin{figure*}[t]
    \centering
    % \vspace{-0.3cm}
    \includegraphics[width=1\textwidth]{figs/short_case.pdf}
%   \vspace{-0.6cm}
    \caption{
    \textbf{Qualitative ablation for STEB and training paradigm.}
     ``TTCA'' and ``MTP'' refer to temporal the text cross attention and mixed training paradigm respectively. The sample without a visual prompt corresponds to the model without spatial image cross attention.
     }
    \label{fig:short_case}
    \vspace{-0.2cm}
\end{figure*}

\begin{table}[t]
    \centering
    \setlength\tabcolsep{6pt}
    \resizebox{0.86\linewidth}{!}{%
    \begin{tabular}{c c c c c c}
        \toprule
        \textbf{TTCA} & \textbf{MTP} & \textbf{0/16 ()} & \textbf{1/15 ()} & \textbf{3/13 ()} & \textbf{5/11 ()}  \\
        \midrule
        \xmark & \xmark & 348.36 & 217.39 & 297.37 & 246.56 \\
        \cmark & \xmark & 333.63 & 178.27 & 230.56 & 193.52 \\
        \cmark & \cmark & \textbf{257.70} & \textbf{123.51} & \textbf{118.02} & \textbf{109.31} \\
        \bottomrule
        \vspace{-1.75em}
    \end{tabular}
    }
    \caption{\textbf{Ablation for temporal text cross attention and mixed training paradigm}. ``TTCA'' and ``MTP'' refer to temporal text cross attention and mixed training paradigm respectively. ``3/13'' denotes predicting the following 13 frames after being given 3 frames of a ground truth video, and we evaluate the FVD between the newly generated 13 frames and the corresponding frames in the ground truth. 
    ``0/16'', ``1/15'', ``5/11'' follow this pattern.
    }
    \vspace{-0.5cm}
    \label{tab:vimeo}
\end{table}



\noindent\textbf{Spatial-Temporal Enhanced Block}.
Fig. \ref{fig:ica} and Tab. \ref{tab:vimeo} evaluate two important operations in our STEB block,
\textit{i.e.}, 
spatial image cross attention and temporal text cross attention.
With spatial image cross attention in Fig. \ref{fig:ica}, 
the CLIP-I is significantly improved as $\beta$ increased, 
while the CLIP-T reached its maximum at $\beta$$=$$0.5$ on COCO2017 \cite{Lin2014MicrosoftCC}.
With temporal text cross attention (TTCA) in Tab. \ref{tab:vimeo}, 
both T2V generation (0/16) and prediction (1/15, 3/13 and 5/11) show a significant improvement on Vimeo11k. 


\noindent\textbf{Mixed Training Paradigm}.
Tab. \ref{tab:vimeo} also presents a model comparison by incorporating our mixed training paradigm or the random mask training method in \cite{chen2023seine, voleti2022mcvd, Hoppe2022DiffusionMF, Fu2022TellMW, Yu2022MAGVITMG, Chang2022MaskGITMG}. 
The results show that the mixed training paradigm is effective in enhancing the model's T2V generation and prediction capabilities.
More details in the supplementary doc.




\subsection{Visualization}

First,
we visualize long video generation,
by comparison with the well-known Phenaki \cite{villegas2023phenaki}.
Note that,
since Phenaki does not have the open-sourced codes,
we alternatively use the demo shown in its official website.
Specifically,
we feed the same story description into our Vlogger.
As shown in Fig. \ref{fig:movie_case}, 
our Vlogger shows superior and more diverse video content,
compared to Phenaki.
Moreover,
our Vlogger sufficiently exhibits the story with a much longer duration (\textit{i.e.}, 7min59s),
according to our LLM-created script.


We further visualize the ablation of T2V video generation and prediction by ShowMaker.
As depicted in Fig. \ref{fig:short_case}, 
ShowMaker had significant improvements in generation and prediction performance,
by incorporating our ShowMaker designs.
Additionally, 
it can leverage visual prompts in the spatial actor cross attention,
for distinguishing the ``Teddy'' concept within text prompts.
